% Part: proof-theory
% Chapter: proof-search
% Section: introduction

\documentclass[../../../include/open-logic-section]{subfiles}

\begin{document}

\olfileid[id]{pt}{ps}{int}

\olsection{Pendahuluan}

Aksioma dan aturan inferensi sistem !!{proof}, seperti \Log{G3c} atau
\Log{N1i}, menentukan pohon sekuen atau !!{formula}s mana yang dihitung
sebagai !!{proof}s yang benar. Jika diberikan pohon sekuen atau !!{formula}s
seperti itu (mungkin dengan informasi tambahan, misalnya aturan mana yang
diterapkan pada setiap langkah, atau label bagi asumsi dan aturan inferensi
yang melepaskannya), definisi aksioma dan aturan memungkinkan kita memeriksa
apakah pohon tersebut merupakan !!{proof} yang benar. Masalah yang berbeda,
dan biasanya lebih sulit, ialah \emph{menemukan} !!a{proof} bagi suatu sekuen
atau bagi suatu !!{formula} tertentu dari asumsi-asumsi yang diberikan.
Itulah masalah \emph{pencarian bukti}.

Terdapat prosedur yang sangat sederhana bagi pencarian bukti. Kita dapat
memandang !!{formula}s sebagai barisan simbol dari alfabet berhingga. Sekuen
dan pohon sekuen atau !!{formula}s juga dapat dipandang sebagai barisan
!!{formula}s (mungkin dengan tanda kurung khusus untuk menunjukkan
percabangan pohon). Aksioma dan aturan memungkinkan kita memeriksa secara
mekanis apakah suatu barisan simbol merepresentasikan !!{proof} yang benar.
Akibatnya, kita dapat secara mekanis menghasilkan \emph{semua} !!{proof}s
yang mungkin dan benar (dengan menghasilkan semua barisan simbol dari
alfabet yang dipakai untuk merepresentasikan !!{proof}s, lalu memeriksa
setiap barisan calon untuk menentukan apakah ia merepresentasikan
!!{proof} yang benar). Prosedur pencarian bukti yang sederhana itu ialah:
hasilkan semua !!{proof}s yang benar sampai kita menemukan !!{proof} bagi
sekuen yang diberikan (atau bagi !!{formula} yang diberikan dengan asumsi
terbuka di antara asumsi-asumsi yang diberikan).

Metode yang lebih baik ialah menggunakan metode naif yang dipakai untuk
menemukan !!a{proof} dengan tangan: mulai dari sekuen akhir, pilih
!!a{formula}, lalu terapkan suatu aturan inferensi secara ``terbalik'' untuk
menghasilkan satu atau beberapa premis yang, jika mempunyai !!{proof}s,
dapat digabungkan dengan inferensi terakhir untuk membentuk !!a{proof} bagi
sekuen yang diberikan. Metode serupa (meskipun lebih rumit) digunakan untuk
menemukan !!a{proof} dalam sistem deduksi alami.

Namun, metode ini belum menjamin keberhasilan: jika Anda memilih aturan yang
keliru, atau !!{formula}s dalam urutan yang keliru, Anda dapat menemui jalan
buntu. Metode itu juga belum ditentukan secara lengkap: pada setiap tahap,
mungkin terdapat lebih dari satu !!{formula} yang dapat dipilih, atau lebih
dari satu aturan yang dapat diterapkan secara terbalik. Suatu
\emph{algoritma pencarian bukti} adalah prosedur yang ditentukan sepenuhnya
dan akan menemukan !!a{proof} jika bukti itu ada.

Beberapa sistem !!{proof} lebih cocok bagi algoritma pencarian bukti yang
sederhana daripada sistem lainnya. Dalam sistem sekuen, !!{main operator}
suatu !!{formula}, beserta sisi tempat formula itu muncul, menentukan aturan
mana yang diterapkan secara terbalik. Misalnya, jika Anda hendak menemukan
!!a{proof} bagi $\Gamma \Sequent \Delta, !A \land !B$ dengan
$!A \land !B$ muncul di sebelah kanan, terapkan \RightR{\land} secara
terbalik dan hasilkan dua premis bersesuaian
$\Gamma \Sequent \Delta, !A$ dan $\Gamma \Sequent \Delta, !B$. Namun, jika
sistem Anda adalah \Log{G1c}, adanya kontraksi membuat persoalan lebih sulit
karena memberi lebih banyak pilihan. Anda juga dapat menerapkan
\RightR{\Contraction} secara terbalik dan menghasilkan premis
$\Gamma \Sequent \Delta, !A \land !B, !A \land !B$. Setelah itu muncul
premis $\Gamma \Sequent \Delta, !A \land !B, !A \land !B, !A \land !B$,
dan seterusnya. Dalam \Log{G2c}, persoalannya bahkan lebih buruk. Menerapkan
aturan \RightR{\land} secara terbalik juga berarti bahwa Anda harus menebak
$\Gamma_1$, $\Gamma_2$ sedemikian sehingga
$\Gamma = \Gamma_1 \cup \Gamma_2$, serta $\Delta_1$, $\Delta_2$ sedemikian
sehingga $\Delta = \Delta_1 \cup \Delta_2$, lalu mencoba premis
$\Gamma_1 \Sequent \Delta_1, !A$ dan
$\Gamma_2 \Sequent \Delta_2, !B$. Tebakan yang keliru dapat membawa Anda
ke jalan buntu, sehingga Anda harus menelusur balik ke awal dan memilih
$\Gamma_1$, $\Gamma_2$, $\Delta_1$, $\Delta_2$ yang berbeda. Jika
!!{formula} tersebut adalah $!A \lor !B$, Anda harus memilih salah satu dari
dua kemungkinan bagi \RightR{\lor}, yang menghasilkan salah satu dari dua
premis $\Gamma \Sequent \Delta, !A$ atau
$\Gamma \Sequent \Delta, !B$. Pilihan yang keliru kembali mengharuskan
penelusuran balik.

Sistem \Log{G3c} tidak mempunyai masalah-masalah ini. Sistem itu tidak
mempunyai aturan struktural, dan pilihan aturan inferensi ditentukan oleh
!!{main operator} serta sisi tempat !!{formula} tersebut muncul. Karena itu,
\Log{G3c} lebih cocok bagi pencarian bukti daripada \Log{G1c} atau
\Log{G2c}. Pencarian yang berhasil menghasilkan !!a{proof} dalam
\Log{G3c}, dan !!{proof} tersebut kemudian dapat diterjemahkan ke \Log{G1c}
atau \Log{G2c} (atau bahkan \Log{N1c}). Jadi, algoritma pencarian bukti bagi
\Log{G3c}, bersama algoritma penerjemahan !!{proof}s~\Log{G3c} menjadi
!!{proof}s dalam sistem lain, juga memecahkan masalah pencarian bukti bagi
sistem-sistem tersebut.

Bagaimana kita mengetahui bahwa suatu algoritma pencarian bukti menjamin
keberhasilan? Kita harus menunjukkan bahwa jika algoritma gagal menemukan
!!a{proof}, tidak mungkin ada !!{proof}. Bagaimanapun, algoritma yang buruk
dapat gagal menemukan bukti walaupun bukti itu ada. Masalah ini tidak
dihadapi algoritma sederhana di atas karena algoritma tersebut menelusuri
semua !!{proof}s yang mungkin. Namun, algoritma pencarian bukti seharusnya
efisien dan mengeluarkan usaha sesedikit mungkin.

Gagasan kunci untuk menunjukkan bahwa algoritma pencarian bukti menjamin
keberhasilan ialah menunjukkan bahwa, jika algoritma gagal, cara kegagalan
itu sendiri memungkinkan kita membuktikan bahwa sekuen tersebut tidak
mungkin mempunyai !!a{proof}. Caranya ialah menunjukkan bahwa sekuen itu
tidak valid. Dalam logika orde pertama klasik, sekuen valid
$\Gamma \Sequent \Delta$ adalah sekuen yang pada setiap !!{structure},
sekurang-kurangnya satu !!{formula} dalam $\Gamma$ salah atau
sekurang-kurangnya satu !!{formula} dalam~$\Delta$ benar. \Log{G3c}
\emph{sahih}, yakni hanya !!{prove}s sekuen valid. Hal ini dibuktikan melalui
induksi pada !!{proof}s: aksioma jelas valid, dan jika premis-premis suatu
aturan valid, kesimpulannya juga valid. Untuk menunjukkan bahwa algoritma
pencarian bukti menjamin keberhasilan, kita menunjukkan bahwa jika pencarian
!!a{proof} bagi $\Gamma \Sequent \Delta$ gagal, kita dapat mendefinisikan
!!a{structure} yang membuat semua !!{formula}s dalam $\Gamma$ benar dan
semua !!{formula}s dalam~$\Delta$ salah. Hal ini juga menetapkan bahwa
\Log{G3c} \emph{lengkap}: jika $\Gamma \Sequent \Delta$ valid, terdapat
!!{proof}~\Log{G3c} baginya. Karena setiap pencarian bukti yang gagal
menunjukkan bahwa $\Gamma \Sequent \Delta$ tidak valid, setiap pencarian
bukti bagi sekuen valid harus berhasil.

\end{document}
